% ส่วนเพิ่มเติมในบทที่ 3 ต่อจากหัวข้อ 3.5 และก่อนบทที่ 4
\section*{3.6 การพัฒนานโยบายการตัดสินใจด้วยขั้นตอนวิธีเชิงพันธุกรรม}
งานในส่วนนี้เป็นการพัฒนาต่อจากกระบวนการเชื่อมโยงอัตลักษณ์ที่อธิบายไว้ในหัวข้อ 3.1--3.5 โดยยังใช้การเตรียมข้อมูล การสร้างคู่ผู้สมัคร และคะแนนจากแบบจำลองเป็นพื้นฐานเดิม สิ่งที่เพิ่มเติมคือวิธีนำคะแนนและหลักฐานของแต่ละคู่มาใช้ตัดสินใจว่า ควรรวมข้อมูลอัตโนมัติ ส่งให้เจ้าหน้าที่ตรวจสอบ หรือไม่รวมข้อมูล ผู้วิจัยจึงเพิ่มขั้นตอนการค้นหานโยบายการตัดสินใจด้วยขั้นตอนวิธีเชิงพันธุกรรม (Genetic Algorithm: GA) เพื่อให้การเลือกเกณฑ์มีหลักฐานจากชุดข้อมูลตรวจสอบรองรับ และสะท้อนภาระงานที่เกิดขึ้นในระบบ CRM ได้ชัดเจนขึ้น

การพัฒนาส่วนนี้อยู่ภายใต้กรอบ CRISP-DM เดิม โดยขยายรายละเอียดในขั้นการสร้างแบบจำลองและการประเมินผล รวมถึงเตรียมผลการตัดสินใจสำหรับกระบวนการปลายทาง ขอบเขตของงานที่เพิ่มจึงครอบคลุมการออกแบบนโยบาย การค้นหาพารามิเตอร์ การทดสอบ และการส่งออกผลรายคู่ ส่วนผลเชิงตัวเลขจากการเปรียบเทียบนโยบายควรนำเสนอในบทผลการทดลองแยกจากวิธีดำเนินงานนี้

\subsection*{3.6.1 เหตุผลและขอบเขตของการพัฒนาเพิ่มเติม}
ในระบบเดิม แบบจำลองให้คะแนนว่าคู่โปรไฟล์มีแนวโน้มเป็นบุคคลเดียวกันมากน้อยเพียงใด จากนั้นจึงใช้เกณฑ์สองระดับแบ่งคำตัดสินเป็น MATCH, REVIEW และ NO\_MATCH วิธีนี้ทำให้ระบบมีทางเลือกสำหรับคู่ที่ยังไม่แน่ชัด อย่างไรก็ตาม การกำหนดเกณฑ์ด้วยค่าคงที่อาจยังไม่สอดคล้องกับสัดส่วนข้อมูล ความผิดพลาด และจำนวนคู่ที่เจ้าหน้าที่ต้องตรวจสอบในแต่ละชุดข้อมูล

ตัวอย่างเช่น การลดเกณฑ์ MATCH อาจทำให้ระบบเชื่อมโยงคู่ที่ถูกต้องได้มากขึ้น แต่ก็อาจรวมคนละคนเข้าด้วยกันมากขึ้นด้วย ในทางกลับกัน การเพิ่มเกณฑ์อาจลดการรวมผิด แต่ทำให้คู่จำนวนมากค้างอยู่ใน REVIEW การเลือกนโยบายจึงต้องพิจารณาทั้งความถูกต้องของการรวมข้อมูลและภาระงานที่ตามมา ผู้วิจัยใช้ GA ค้นหาชุดพารามิเตอร์ที่ทำให้ต้นทุนการตัดสินใจบนชุด validation ต่ำลง แล้วจึงตรวจสอบผลด้วยตัวชี้วัดด้านความแม่นยำและปริมาณงานร่วมกัน

ในการเปรียบเทียบนโยบาย ผู้วิจัยตรึงคะแนนจากแบบจำลองไว้ เพื่อให้ความแตกต่างของผลลัพธ์อธิบายได้จากการเปลี่ยนชั้นการตัดสินใจ ส่วนการปรับคุณลักษณะหรือฝึกแบบจำลองในสายการทดลองอื่นถือเป็นงานอีกส่วนหนึ่ง น้ำหนักที่ GA ค้นหาในขั้นตอนนี้จึงเป็นน้ำหนักของนโยบายการตัดสินใจ ไม่ใช่น้ำหนักภายในแบบจำลอง Gradient Boosting และไม่ใช่ค่าความสำคัญของคุณลักษณะจากแบบจำลองโดยตรง

\subsection*{3.6.2 ลำดับการพัฒนาจากกฎเดิมสู่นโยบายแบบถ่วงน้ำหนัก}
ผู้วิจัยพัฒนานโยบายเป็นลำดับ โดยเริ่มจากกฎที่ใช้คะแนนจากแบบจำลองและเกณฑ์สองระดับเป็นฐานเปรียบเทียบ จากนั้นเพิ่มการใช้ความคล้ายของชื่อช่วยตัดสินคู่ที่อยู่ในช่วง REVIEW และจึงขยายไปสู่นโยบายที่รวมหลักฐานหลายด้านเป็นคะแนนเดียว ตารางที่ 3.18 สรุปความแตกต่างของแนวทางที่ใช้ในการศึกษา

\begin{center}\small
\begin{tabularx}{\linewidth}{p{3.4cm}XX}
\toprule
แนวทาง & สิ่งที่ใช้ตัดสินใจ & บทบาทในการทดลอง \\
\midrule
เกณฑ์สองระดับจากคะแนนโมเดล & คะแนนโมเดลและเกณฑ์ MATCH กับ REVIEW & ใช้ตรวจว่าการปรับเกณฑ์เพียงอย่างเดียวเพียงพอหรือไม่ \\
GA ที่ใช้ชื่อช่วยตัดสิน & คะแนนโมเดล เกณฑ์สองระดับ และเกณฑ์ความคล้ายของชื่ออีกสองค่า & เป็นแนวทางต่อยอดระยะแรกและฐานเปรียบเทียบของนโยบายใหม่ \\
GA แบบถ่วงน้ำหนักขั้นเดียว & คะแนนโมเดลร่วมกับหลักฐานที่มีข้อมูลของแต่ละคู่ และเกณฑ์สองระดับ & ค้นหาน้ำหนักและเกณฑ์ร่วมกัน แล้วตัดสินจากคะแนนรวมครั้งเดียว \\
\bottomrule
\end{tabularx}
\par ตารางที่ 3.18 ลำดับการพัฒนานโยบายการตัดสินใจ
\end{center}

นโยบายที่ใช้ชื่อช่วยตัดสินมีโครโมโซม 4 ค่า ได้แก่ เกณฑ์ MATCH เกณฑ์ REVIEW เกณฑ์ความคล้ายของชื่อสำหรับเลื่อนคู่ขึ้นเป็น MATCH และเกณฑ์สำหรับลดคู่ลงเป็น NO\_MATCH โดยเกณฑ์ของชื่อทำงานเฉพาะคู่ที่คะแนนโมเดลอยู่ในช่วง REVIEW เท่านั้น ข้อจำกัดของแนวทางนี้คือ ชื่อถูกกำหนดให้เป็นหลักฐานหลักตั้งแต่ก่อนเริ่มค้นหา GA จึงเลือกค่าจุดตัดของชื่อได้ แต่ยังไม่ได้เปรียบเทียบว่าหลักฐานชนิดอื่นเหมาะสมกว่าหรือไม่

เพื่อศึกษาข้อจำกัดดังกล่าว ผู้วิจัยเพิ่มการทดลองใช้หลักฐานแต่ละกลุ่มเป็นตัวช่วยตัดสิน และต่อยอดเป็นนโยบายแบบถ่วงน้ำหนักขั้นเดียว (One-Stage Weighted Policy) ซึ่งรวมคะแนนโมเดลและหลักฐานที่มีอยู่ก่อนเปรียบเทียบกับเกณฑ์ MATCH และ REVIEW แนวทางนี้ไม่มีเกณฑ์เลื่อนขึ้นและลดลงของชื่อซ้อนอยู่ในช่วง REVIEW อีกชั้นหนึ่ง แต่ยังเปิดให้ GA เลือกน้ำหนักของชื่อได้ตามผลบน validation

\subsection*{3.6.3 การเตรียมข้อมูลสำหรับชั้นการตัดสินใจ}
ข้อมูลสำหรับนโยบายต้องอ้างอิงคู่เดียวกับที่ระบบให้คะแนนไว้จริง ผู้วิจัยจึงจัดเตรียมตารางคุณลักษณะสำหรับการตัดสินใจ โดยเก็บรหัสโปรไฟล์ทั้งสองฝั่ง คะแนนโมเดล กลุ่มข้อมูลที่ใช้ประเมิน ค่าความจริงสำหรับคำนวณผล และหลักฐานแต่ละด้านไว้ด้วยกัน พร้อมตัวบ่งชี้ว่าหลักฐานด้านนั้นใช้งานได้หรือไม่ การมีรหัสคู่ช่วยตรวจสอบว่าคะแนนและคุณลักษณะไม่ได้ถูกนำไปจับกับคู่ผิดระหว่างการประมวลผลแบบแบ่งส่วน

โครงสร้างนโยบายกำหนดช่องข้อมูลไว้ 8 ช่อง ได้แก่ คะแนนโมเดล ชื่อ ประวัติย่อ ลิงก์ ตำแหน่งที่อยู่ สัญญาณทางสังคม ภาพ และรูปแบบการเขียน ดังตารางที่ 3.19 ค่าหลักฐานถูกจัดให้อยู่ในช่วง 0 ถึง 1 และมีทิศทางเดียวกัน คือค่าสูงแสดงความสอดคล้องมากกว่า ส่วนคุณลักษณะที่เดิมเป็นค่าความแตกต่างต้องแปลงทิศทางก่อนนำมาใช้

\begin{center}\small
\begin{tabularx}{\linewidth}{p{2.5cm}X}
\toprule
กลุ่มข้อมูล & ความหมายในชั้นการตัดสินใจ \\
\midrule
คะแนนโมเดล & คะแนนความเป็นบุคคลเดียวกันจากแบบจำลองที่ตรึงไว้ \\
ชื่อ & ความคล้ายของชื่อผู้ใช้หรือชื่อที่แสดงในโปรไฟล์ \\
ประวัติย่อ & ความคล้ายของข้อความหรือความหมายในประวัติย่อ \\
ลิงก์ & ความสอดคล้องของ URL หรือโดเมนที่อ้างอิง \\
ตำแหน่งที่อยู่ & ความคล้ายของข้อมูลสถานที่เมื่อมีข้อมูลที่ใช้งานได้ \\
สัญญาณทางสังคม & ความสอดคล้องของการกล่าวถึงบัญชีหรือแฮชแท็กในคุณลักษณะที่เตรียมไว้ \\
ภาพ & หลักฐานจากภาพหรือข้อความประกอบภาพตามที่มีในตารางคุณลักษณะ \\
รูปแบบการเขียน & ความสอดคล้องของลักษณะการเขียน เช่น ความยาวและการใช้เครื่องหมาย \\
\bottomrule
\end{tabularx}
\par ตารางที่ 3.19 กลุ่มข้อมูลที่โครงสร้างนโยบายรองรับ
\end{center}

การรองรับ 8 ช่องไม่ได้หมายความว่าทุกรอบทดลองใช้ข้อมูลได้ครบทั้ง 8 ด้าน ในรอบข้อมูลเต็มวันที่ 14 กันยายน 2569 ตารางคุณลักษณะไม่มีช่อง social และไม่มีรายการ location ที่พร้อมใช้งาน จึงปิดการใช้งานสองด้านนี้และกำหนดน้ำหนักเป็นศูนย์ตั้งแต่เริ่มค้นหา ส่วนกลุ่มที่มีข้อมูลให้ใช้เปิดให้ GA ค้นหาน้ำหนักได้ตามการตั้งค่าของรอบนั้น การเปิดใช้งานในรอบนี้พิจารณาว่ามีหลักฐานด้านนั้นปรากฏใน validation หรือไม่ โดยไม่ได้กำหนดเกณฑ์ขั้นต่ำด้านสัดส่วนความครอบคลุมและจำนวนคู่บวกเพิ่มเติม

ผู้วิจัยตรวจสอบสัดส่วนคู่ที่มีหลักฐานพร้อมใช้งานควบคู่กับน้ำหนักที่ได้ เพื่อแยกกรณี ``ไม่มีข้อมูลให้ใช้'' ออกจากกรณี ``มีข้อมูล แต่ GA เลือกน้ำหนักเป็นศูนย์'' นอกจากนี้ หลักฐานกลุ่มภาพอาจรวมสัญญาณจากคำบรรยายภาพด้วย สัดส่วนความพร้อมของกลุ่มภาพจึงไม่ควรตีความเป็นสัดส่วนคู่ที่มีภาพจริงครบทั้งสองฝั่งโดยอัตโนมัติ

\subsection*{3.6.4 การคำนวณคะแนนเมื่อข้อมูลแต่ละคู่มีไม่ครบ}
โปรไฟล์แต่ละคู่มีข้อมูลไม่เท่ากัน หากแทนข้อมูลที่ขาดด้วยศูนย์แล้วเฉลี่ยทุกด้านเท่ากัน คู่ที่ไม่มีประวัติย่อหรือภาพอาจได้คะแนนต่ำลงเพียงเพราะไม่มีข้อมูล ผู้วิจัยจึงแยกค่าหลักฐานออกจากตัวบ่งชี้ความพร้อมใช้งาน และปรับสัดส่วนน้ำหนักตามข้อมูลที่มีจริงในแต่ละคู่

กำหนดให้ $s_{ik}$ เป็นค่าหลักฐานช่องที่ $k$ ของคู่ที่ $i$, $m_{ik}$ เป็นตัวบ่งชี้ความพร้อมใช้งาน โดยมีค่า 1 เมื่อใช้ได้และ 0 เมื่อใช้ไม่ได้ และ $w_k$ เป็นน้ำหนักที่ GA ค้นหา คะแนนรวมของแต่ละคู่คำนวณได้ดังนี้
\begin{equation*}
S_i=\frac{\sum_{k=1}^{8} w_k m_{ik}s_{ik}}{\sum_{k=1}^{8} w_k m_{ik}},\qquad w_k\geq 0,\quad \sum_{k=1}^{8}w_k=1
\end{equation*}

สูตรนี้ใช้เมื่อผลรวมน้ำหนักของช่องที่พร้อมใช้งานมากกว่าศูนย์ ตัวอย่างเช่น หากนโยบายให้น้ำหนักคะแนนโมเดล 0.2 ชื่อ 0.5 และประวัติย่อ 0.3 แต่คู่หนึ่งไม่มีประวัติย่อ ระบบจะหารด้วย 0.7 ซึ่งเป็นน้ำหนักรวมของสองด้านที่ยังใช้ได้ แทนการหารด้วย 1 จึงไม่ลดคะแนนเพราะการขาดประวัติย่อเพียงอย่างเดียว หากไม่มีหลักฐานที่มีน้ำหนักให้ใช้เลย ระบบกำหนดคู่นั้นเป็น REVIEW

จากนั้นระบบเปรียบเทียบคะแนนรวมกับเกณฑ์สองระดับเพียงครั้งเดียว ดังนี้
\begin{equation*}
D_i=\begin{cases}
\mathrm{MATCH}, & S_i\geq t_{\mathrm{match}},\\
\mathrm{NO\_MATCH}, & S_i<t_{\mathrm{review}},\\
\mathrm{REVIEW}, & t_{\mathrm{review}}\leq S_i<t_{\mathrm{match}}.
\end{cases}
\end{equation*}

MATCH หมายถึงคู่ที่ผ่านเกณฑ์สำหรับการเชื่อมโยงอัตโนมัติ REVIEW หมายถึงคู่ที่ยังต้องตรวจสอบเพิ่มเติม และ NO\_MATCH หมายถึงคู่ที่ไม่ผ่านเกณฑ์การเชื่อมโยงในรอบนั้น คะแนน $S_i$ เป็นคะแนนรวมของนโยบาย จึงไม่ควรเรียกว่าเป็นความน่าจะเป็นที่ผ่านการปรับเทียบแล้ว เว้นแต่มีการประเมินการปรับเทียบของคะแนนรวมนี้เพิ่มเติม

\subsection*{3.6.5 การกำหนดโครโมโซมและเงื่อนไขการค้นหา}
โครโมโซมของนโยบายแบบถ่วงน้ำหนักประกอบด้วย 10 ค่า คือเกณฑ์สองระดับและน้ำหนักอีก 8 ค่า เรียงตามลำดับดังนี้
\begin{equation*}
\boldsymbol{g}=[t_m,t_r,w_p,w_n,w_b,w_u,w_l,w_s,w_i,w_y]
\end{equation*}
โดย $t_m$ และ $t_r$ แทนเกณฑ์ MATCH และ REVIEW ส่วนน้ำหนัก $w_p$ ถึง $w_y$ แทนคะแนนโมเดล ชื่อ ประวัติย่อ ลิงก์ ตำแหน่งที่อยู่ สัญญาณทางสังคม ภาพ และรูปแบบการเขียนตามลำดับ

การค้นหาในโค้ดกำหนดให้ $0.50\leq t_m\leq0.999$ และ $0.05\leq t_r\leq t_m$ น้ำหนักทุกค่าต้องไม่ติดลบ และปรับผลรวมให้เท่ากับ 1 หลังสร้างหรือเปลี่ยนแปลงโครโมโซม กลุ่มข้อมูลที่ไม่ได้เปิดใช้งานมีน้ำหนักเป็นศูนย์ตลอดการค้นหา เงื่อนไขเหล่านี้ช่วยให้โครโมโซมทุกชุดนำไปตัดสินใจได้จริงและคงลำดับของเกณฑ์สองระดับไว้เสมอ

\subsection*{3.6.6 การกำหนดต้นทุนสำหรับประเมินนโยบาย}
ผู้วิจัยประเมินนโยบายด้วยต้นทุนสามประเภท ได้แก่ การรวมโปรไฟล์ของคนละคนเข้าด้วยกัน การไม่รวมคู่ที่เป็นคนเดียวกัน และการส่งคู่ให้เจ้าหน้าที่ตรวจสอบ โดยใช้สมมติฐานค่าแรง 100 บาทต่อชั่วโมง เวลาตรวจสอบคู่ละ 1 นาที เวลาแก้ไขการรวมผิด 20 นาที และเวลาจัดการคู่ที่ตกหล่น 10 นาที สมมติฐานเหล่านี้เป็นค่าที่กำหนดสำหรับการทดลอง ยังไม่ใช่ต้นทุนที่วัดจากการใช้งาน CRM จริง

เมื่อคิดเฉพาะค่าแรง จะได้ต้นทุนประมาณ 33.33 บาทต่อการรวมผิดหนึ่งคู่ 16.67 บาทต่อคู่ที่ตกหล่นหนึ่งคู่ และ 1.67 บาทต่อการตรวจสอบหนึ่งคู่ เมื่อนำต้นทุนคู่ที่ตกหล่นมาเป็นหน่วยอ้างอิง จะได้อัตราส่วน 2 ต่อ 1 ต่อ 0.1 และกำหนดฟังก์ชันต้นทุนดังนี้
\begin{equation*}
C(\boldsymbol{g})=2\,FP+FN+0.1\,R
\end{equation*}

ในสมการนี้ $FP$ คือจำนวนคู่คนละคนที่ถูกตัดสินเป็น MATCH, $FN$ คือจำนวนคู่คนเดียวกันที่ถูกตัดสินเป็น NO\_MATCH และ $R$ คือจำนวนคู่ทั้งหมดที่ถูกส่งเป็น REVIEW การนับ $FN$ ในต้นทุนจึงหมายถึงการปฏิเสธคู่ที่ควรเชื่อมโยงโดยตรง ส่วนคู่บวกที่อยู่ใน REVIEW ถูกคิดเป็นต้นทุนการตรวจสอบ และไม่ได้ถูกคิดซ้ำเป็น $FN$ ในสมการเดียวกัน

GA ค้นหานโยบายที่ทำให้ต้นทุนดังกล่าวต่ำที่สุดบน validation การลดต้นทุนอาจมาพร้อมกับการเปลี่ยนแปลงของ precision, recall และจำนวน REVIEW จึงต้องรายงานตัวชี้วัดเหล่านี้ร่วมกัน ไม่สามารถสรุปจากต้นทุนเพียงค่าเดียวว่านโยบายดีกว่าในทุกด้าน และไม่ควรแปลงผลต้นทุนเป็นเงินที่องค์กรประหยัดได้จริงโดยไม่มีข้อมูลการปฏิบัติงานรองรับ

\subsection*{3.6.7 ขั้นตอนการค้นหานโยบายด้วย GA}
การค้นหาเริ่มจากสร้างประชากรของชุดพารามิเตอร์หลายชุด แต่ละชุดแทนนโยบายที่เป็นไปได้หนึ่งนโยบาย ผู้วิจัยสุ่มค่าเกณฑ์และน้ำหนักเริ่มต้น แล้วปรับให้อยู่ภายใต้เงื่อนไขที่กำหนด จากนั้นนำทุกนโยบายไปตัดสินคู่ใน validation และคำนวณต้นทุนของแต่ละชุด

ในแต่ละรุ่น ระบบเรียงประชากรตามต้นทุนจากน้อยไปมาก สำหรับตัวค้นหานโยบายแบบถ่วงน้ำหนัก หากต้นทุนเท่ากันจะพิจารณา precision ที่สูงกว่า จำนวน FP ที่น้อยกว่า recall ที่สูงกว่า และจำนวน REVIEW ที่น้อยกว่าตามลำดับ แล้วเก็บนโยบายอันดับต้นไว้เป็นกลุ่มชั้นนำเพื่อใช้สร้างรุ่นถัดไป

การสร้างลูกใช้การเลือกยีนแต่ละตำแหน่งจากพ่อหรือแม่ด้วยโอกาสเท่ากัน จากนั้นกลายพันธุ์แต่ละยีนด้วยความน่าจะเป็น 0.30 โดยเพิ่มค่ารบกวนจากการแจกแจงปกติ ส่วนเบี่ยงเบนมาตรฐานสำหรับเกณฑ์เท่ากับ 0.025 และสำหรับน้ำหนักเท่ากับ 0.05 หลังกลายพันธุ์ ระบบตรวจขอบเขตของเกณฑ์ ปิดน้ำหนักของกลุ่มที่ไม่ใช้งาน และปรับผลรวมน้ำหนักอีกครั้ง กระบวนการนี้ทำซ้ำตามจำนวนรุ่นที่กำหนด

การทดลองบนข้อมูลเต็มใช้ประชากร 64 ชุด จำนวน 60 รุ่น เก็บประชากรชั้นนำ 8 ชุด และทำซ้ำด้วยค่า seed 42, 43 และ 44 ทั้งนโยบายคะแนนอย่างเดียว นโยบายใช้ชื่อช่วยตัดสิน และนโยบายแบบถ่วงน้ำหนักใช้ขนาดประชากรและจำนวนรุ่นเท่ากันในรอบนี้ หลังจบการค้นหา ผู้วิจัยเลือกผลของแต่ละนโยบายข้าม seed โดยเรียงตามต้นทุน FP, FN, REVIEW และ seed ตามลำดับ พร้อมเก็บประวัติการค้นหาไว้ตรวจสอบความเปลี่ยนแปลงของคำตอบ

การใช้หลาย seed ช่วยตรวจสอบผลจากการสุ่มเริ่มต้น แต่ไม่ได้รับประกันว่าพบคำตอบที่ดีที่สุดของพื้นที่ค้นหาทั้งหมด นอกจากนี้ ค่าตั้งต้นของการทดลองย่อยอาจต่างจากรอบข้อมูลเต็ม จึงต้องอ้างอิงการตั้งค่าจากไฟล์ของรอบที่นำมารายงานเสมอ

\section*{3.7 การออกแบบการประเมินนโยบายที่พัฒนาเพิ่มเติม}
\subsection*{3.7.1 การแบ่งบทบาทของข้อมูลและการตรึงนโยบาย}
ผู้วิจัยแยกบทบาทของการฝึกแบบจำลอง การเลือกนโยบาย และการทดสอบออกจากกัน โดยใช้คะแนนและการแบ่งกลุ่มที่บันทึกไว้ในชุดข้อมูลของแต่ละการทดลอง การค้นหาพารามิเตอร์ของ GA ใช้เฉพาะคู่ validation ส่วน test ใช้ประเมินนโยบายหลังเลือกเสร็จแล้ว การเปรียบเทียบแต่ละนโยบายจึงใช้คู่ข้อมูล คะแนนโมเดล และนิยามต้นทุนเดียวกัน

สำหรับรอบข้อมูลเต็ม มีคู่ทั้งหมด 2,086,245 คู่ แบ่งเป็นคู่ที่ใช้คะแนนโมเดล 2,073,842 คู่ และคู่ในชั้น Exact-First 12,403 คู่ โดย GA ใช้คู่ validation ที่ผ่านการให้คะแนนโมเดลจำนวน 1,010,926 คู่ในการค้นหา หลังเลือกนโยบาย ผู้วิจัยบันทึกค่าเกณฑ์ น้ำหนัก seed รายชื่อกลุ่มข้อมูลที่เปิดใช้งาน และรหัสตรวจสอบไฟล์ต้นทางไว้ก่อนเรียกขั้นตอนประเมิน test

การตรึงนโยบายก่อนทดสอบช่วยป้องกันการนำผล test ในรอบนั้นกลับไปปรับพารามิเตอร์ อย่างไรก็ตาม การมีไฟล์ตรึงนโยบายไม่ได้ยืนยันด้วยตัวเองว่าคะแนนต้นทางทุกชุดถูกสร้างโดยปราศจากข้อมูลรั่วไหล หรือว่าชุดทดสอบไม่เคยถูกใช้ในงานก่อนหน้า การกล่าวถึงการแยกข้อมูลระดับบุคคลหรือการทดสอบกับชุดใหม่จึงต้องอาศัยหลักฐานของกระบวนการต้นทางเพิ่มเติมด้วย

\subsection*{3.7.2 ฐานเปรียบเทียบและขอบเขตการเปลี่ยนแปลง}
การประเมินใช้สี่แนวทาง ได้แก่ ผลการตัดสินใจจากระบบเดิม นโยบายคะแนนโมเดลอย่างเดียวที่ค้นหาเกณฑ์สองระดับด้วย GA นโยบายที่ใช้ชื่อช่วยตัดสินด้วยพารามิเตอร์สี่ค่า และนโยบายแบบถ่วงน้ำหนักสิบค่า การมีฐานคะแนนอย่างเดียวช่วยแยกผลของการปรับเกณฑ์ออกจากผลของการเพิ่มหลักฐานหลายด้าน ส่วนนโยบายใช้ชื่อช่วยตัดสินทำให้เห็นผลของการเปลี่ยนรูปแบบกฎจากการตัดสินในช่วง REVIEW มาเป็นคะแนนรวมขั้นเดียว

คู่ที่ผ่านชั้น Exact-First ยังคงคำตัดสิน MATCH ตามระบบเดิมทุกแนวทาง จึงไม่ใช่ส่วนที่ GA ปรับปรุง ความผิดพลาดของชั้นดังกล่าวยังต้องนับเมื่อรายงานผลรวม ขณะเดียวกัน ผลบนคู่ที่ใช้คะแนนโมเดลเพียงอย่างเดียวต้องแยกจากผลที่รวม Exact-First แล้ว เพื่อไม่ให้ความแม่นยำของกฎเดิมถูกตีความเป็นผลของ GA ทั้งหมด

\subsection*{3.7.3 ตัวชี้วัดและการอ่านผลลัพธ์สามทาง}
ผู้วิจัยรายงานต้นทุน จำนวน FP จำนวน FN จำนวน REVIEW, precision, recall และ F1 ร่วมกัน โดย precision ของการรวมอัตโนมัติคำนวณจากคู่บวกที่เป็น MATCH หารด้วยคู่ทั้งหมดที่เป็น MATCH ส่วน recall คำนวณจากคู่บวกที่เป็น MATCH หารด้วยคู่บวกทั้งหมดในขอบเขตที่กำลังประเมิน

\begin{equation*}
\mathrm{Precision}=\frac{TP}{TP+FP},\qquad
\mathrm{Recall}=\frac{TP}{N_+},\qquad
F1=\frac{2\,\mathrm{Precision}\,\mathrm{Recall}}{\mathrm{Precision}+\mathrm{Recall}}
\end{equation*}

กำหนดให้ $TP$ เป็นคู่บวกที่ตัดสินเป็น MATCH และ $N_+$ เป็นคู่บวกทั้งหมดในชุดประเมิน คู่บวกที่อยู่ใน REVIEW ยังไม่ถือว่ารวมสำเร็จ จึงยังอยู่ในตัวหารของ recall ด้วย ดังนั้น ค่า FN ที่แสดงในฟังก์ชันต้นทุนไม่สามารถนำไปแทนตัวหารของ recall เป็น $TP+FN$ ได้ทันที เพราะจะขาดคู่บวกที่รอตรวจสอบ หากตัวหารเป็นศูนย์ต้องระบุว่าไม่มีตัวอย่างรองรับการตีความ แทนการมองค่าที่โปรแกรมเติมให้เป็นหลักฐานของประสิทธิภาพ

การรายงานแยกผลบน validation ผลบน test และผลการนำไปใช้กับคู่ทั้งหมด โดยผลรวมทุกคู่ใช้แสดงปริมาณงานที่นโยบายก่อให้เกิดขึ้นในข้อมูลชุดนั้น และไม่ใช้แทนผลทดสอบบนข้อมูลที่กันไว้ กลุ่มที่บันทึกเป็น drop ต้องแสดงแยกตามสถานะของข้อมูลต้นทาง และไม่เรียกว่าเป็นชุดทดสอบเพิ่มเติม

นอกจากนี้ การประเมินบนคู่ผู้สมัครยังไม่ครอบคลุมคู่จริงที่หลุดไปตั้งแต่ขั้นสร้างคู่ ผู้วิจัยจึงอ่านผลนโยบายร่วมกับความครอบคลุมของ candidate generation ที่อธิบายไว้ในหัวข้อ 3.5 การปรับ GA ไม่สามารถกู้คืนคู่ที่ไม่เคยถูกส่งเข้ามาให้ตัดสินได้ หากไม่ได้ปรับขั้นสร้างคู่ด้วย

\subsection*{3.7.4 การตรวจสอบน้ำหนักและข้อมูลที่ขาดหาย}
ผู้วิจัยรายงานน้ำหนักที่ตรึงไว้พร้อมสัดส่วนความพร้อมใช้งานของหลักฐาน เพื่อช่วยอธิบายว่าระบบอาศัยข้อมูลใดในการตัดสินใจ น้ำหนักรวมของนโยบายยังต้องแยกจากน้ำหนักที่มีผลจริงต่อแต่ละคู่ เพราะสูตรคำนวณจะปรับสัดส่วนใหม่เมื่อบางด้านไม่มีข้อมูล ตัวอย่างเช่น น้ำหนักภาพในนโยบายอาจมีค่าน้อย แต่ในคู่ที่ไม่มีหลักฐานชนิดอื่น สัดส่วนที่มีผลจริงอาจแตกต่างจากน้ำหนักรวมมาก

การวิเคราะห์จึงพิจารณาทั้งน้ำหนัก ความครอบคลุม ความเปลี่ยนแปลงระหว่าง seed และข้อผิดพลาดของคู่ที่มีหรือไม่มีหลักฐานแต่ละด้าน การอ้างว่ากลุ่มใดสำคัญที่สุดต่อแบบจำลองยังต้องใช้การวิเคราะห์แบบจำลองโดยตรงเพิ่มเติม ไม่สามารถสรุปจากน้ำหนักนโยบายเพียงอย่างเดียวได้

\section*{3.8 ขั้นตอนการดำเนินงานส่วนที่พัฒนาเพิ่มเติม}
เพื่อให้เห็นความเชื่อมโยงกับงานเดิม ผู้วิจัยจัดลำดับงานเพิ่มเติมตั้งแต่ตรวจสอบข้อมูลต้นทางจนถึงจัดเตรียมผลสำหรับระบบ CRM เป็น 9 ขั้นตอนดังนี้

\subsection*{ขั้นตอนที่ 1 ทบทวนผลและข้อจำกัดของระบบเดิม}
รวบรวมคะแนนรายคู่ คำตัดสินเดิม และผลประเมินของกระบวนการ Exact-First ร่วมกับเกณฑ์สองระดับ ตรวจสอบจำนวนคู่ที่รวมผิด คู่ที่ปฏิเสธผิด และคู่ที่รอตรวจสอบ เพื่อกำหนดว่าการพัฒนาครั้งนี้จะปรับวิธีตัดสินใจในส่วนใด ผลที่ต้องได้คือฐานเปรียบเทียบที่มีขอบเขตและนิยามตัวชี้วัดตรงกัน

\subsection*{ขั้นตอนที่ 2 จัดเตรียมคู่ข้อมูลและตรวจความถูกต้อง}
แยกคู่ Exact-First ออกจากคู่ที่ใช้คะแนนโมเดล ตรวจรหัสคู่ซ้ำ จำนวนแถว และรหัสตรวจสอบไฟล์ต้นทาง จากนั้นจัดคู่ validation และ test ตามข้อมูลการแบ่งกลุ่มที่บันทึกไว้ ขั้นตอนนี้ช่วยให้เปรียบเทียบทุกนโยบายกับประชากรเดียวกัน และติดตามได้ว่าผลที่รายงานมาจากข้อมูลรุ่นใด

\subsection*{ขั้นตอนที่ 3 สร้างหลักฐานสำหรับการตัดสินใจ}
เชื่อมข้อมูลคู่กับคุณลักษณะชื่อ ประวัติย่อ ลิงก์ ตำแหน่งที่อยู่ สัญญาณทางสังคม ภาพ และรูปแบบการเขียนตามที่มีจริง จัดทิศทางคะแนนให้สอดคล้องกัน และสร้างตัวบ่งชี้ความพร้อมใช้งานแยกจากค่าคะแนน สำหรับข้อมูลจำนวนมากให้ประมวลผลเป็นส่วนย่อย และตรวจรหัสคู่ทุกส่วนก่อนรวมผล

\subsection*{ขั้นตอนที่ 4 ตรวจความครอบคลุมและกำหนดกลุ่มที่ใช้งาน}
นับจำนวนและสัดส่วนคู่ที่ใช้หลักฐานแต่ละด้านได้ ตรวจกรณีข้อมูลหายหรือไม่มีช่องข้อมูล และกำหนดรายชื่อกลุ่มที่เปิดให้ GA ค้นหาน้ำหนักจาก validation สำหรับรอบข้อมูลเต็ม ปิด location และ social ตามสภาพข้อมูลจริง พร้อมบันทึกเหตุผล เพื่อไม่ตีความการไม่มีข้อมูลเป็นผลเชิงลบของหลักฐานด้านนั้น

\subsection*{ขั้นตอนที่ 5 กำหนดนโยบายเปรียบเทียบและต้นทุน}
เตรียมนโยบายคะแนนอย่างเดียว นโยบายใช้ชื่อช่วยตัดสิน และนโยบายแบบถ่วงน้ำหนัก กำหนดต้นทุนเป็น $2FP+FN+0.1R$ เหมือนกันทุกแนวทาง และระบุให้ชัดว่าคู่บวกใน REVIEW คิดเป็นงานตรวจสอบ การกำหนดวิธีนับก่อนเริ่มค้นหาช่วยลดความคลาดเคลื่อนในการเปรียบเทียบผล

\subsection*{ขั้นตอนที่ 6 ค้นหาพารามิเตอร์บน validation}
ใช้ GA ค้นหาค่าเกณฑ์และน้ำหนักตามโครงสร้างของแต่ละนโยบาย ทำซ้ำด้วย seed ที่กำหนด และบันทึกต้นทุนกับพารามิเตอร์ระหว่างการค้นหา เปรียบเทียบผลจากการตั้งค่าของรอบเดียวกัน แล้วเลือกชุดที่ดีที่สุดตามลำดับเกณฑ์ที่กำหนดไว้ โดยยังไม่นำผล test มาใช้เลือก

\subsection*{ขั้นตอนที่ 7 บันทึกและตรึงนโยบายก่อนทดสอบ}
จัดเก็บเกณฑ์ น้ำหนัก seed การตั้งค่า GA กลุ่มข้อมูลที่เปิดใช้งาน และรหัสตรวจสอบไฟล์ต้นทางลงในไฟล์นโยบายที่ตรึงไว้ ขั้นตอนนี้ทำให้ย้อนตรวจได้ว่าชุดพารามิเตอร์ที่ใช้ประเมิน test เป็นชุดเดียวกับที่เลือกจาก validation และสามารถนำกลับมาใช้ตัดสินข้อมูลได้ซ้ำ

\subsection*{ขั้นตอนที่ 8 ประเมิน test และนำไปตัดสินคู่ทั้งหมด}
นำพารามิเตอร์ที่ตรึงไว้ไปประเมิน test โดยไม่ปรับค่าเพิ่มเติม จากนั้นคำนวณคำตัดสินสำหรับคู่ทั้งหมดโดยคง Exact-First ตามเดิม และส่งออกผลแยกตามนโยบาย รายงานผล validation, test, drop และภาพรวมแยกกัน เพื่อให้เห็นทั้งคุณภาพในขอบเขตทดสอบและปริมาณงานที่เกิดเมื่อใช้กับข้อมูลทั้งหมด

\subsection*{ขั้นตอนที่ 9 วิเคราะห์ข้อผิดพลาดและเตรียมผลส่งต่อ CRM}
ตรวจคู่ที่ตัดสินต่างจากระบบเดิม โดยเฉพาะคู่ที่เปลี่ยนจาก REVIEW เป็น MATCH หรือ NO\_MATCH เพื่ออธิบายว่าภาระตรวจสอบที่ลดลงสัมพันธ์กับความผิดพลาดอย่างไร จัดเก็บรหัสคู่ คะแนน คำตัดสิน และรุ่นนโยบายสำหรับส่งต่อกระบวนการรวมโปรไฟล์หรือคิวตรวจสอบ ผลส่งออกนี้เป็นหลักฐานของการประมวลผลในโครงงาน ยังไม่ใช่หลักฐานว่าเจ้าหน้าที่ได้ตรวจครบทุกคู่ หรือระบบ CRM ใช้งานจริงและเกิดผลประหยัดตามสมมติฐานแล้ว

\begin{center}\small
\begin{tabularx}{\linewidth}{p{3.4cm}X}
\toprule
ช่วงดำเนินงาน & ผลงานที่จัดเก็บเพื่อตรวจสอบย้อนกลับ \\
\midrule
เตรียมข้อมูล & ตารางหลักฐานรายคู่ ตัวบ่งชี้ความพร้อมใช้งาน และข้อมูลตรวจสอบแหล่งที่มา \\
ค้นหานโยบาย & ผล validation ของแต่ละ seed และประวัติการค้นหา GA \\
ตรึงนโยบาย & เกณฑ์ น้ำหนัก การตั้งค่า และไฟล์นโยบายก่อนประเมิน test \\
ประเมินและส่งออก & ตารางตัวชี้วัดแยกกลุ่มข้อมูล และคะแนนกับคำตัดสินรายคู่ของแต่ละนโยบาย \\
\bottomrule
\end{tabularx}
\par ตารางที่ 3.20 ผลงานจากขั้นตอนการพัฒนานโยบายเพิ่มเติม
\end{center}

การดำเนินงานส่วนนี้ทำให้กระบวนการเดิมมีขั้นตอนเลือกนโยบายที่ตรวจสอบย้อนกลับได้ ตั้งแต่หลักฐานที่ใช้ ค่าพารามิเตอร์ที่เลือก ไปจนถึงผลกระทบต่อการรวมข้อมูลและคิวตรวจสอบ โดยยังรักษาความเชื่อมโยงกับการเตรียมข้อมูล การสร้างคู่ และการให้คะแนนที่ได้พัฒนาไว้ก่อนหน้า
